% Options for packages loaded elsewhere
\PassOptionsToPackage{unicode}{hyperref}
\PassOptionsToPackage{hyphens}{url}
\documentclass[
  english,
]{article}
\usepackage{xcolor}
\usepackage{amsmath,amssymb}
\setcounter{secnumdepth}{-\maxdimen} % remove section numbering
\usepackage{iftex}
\ifPDFTeX
  \usepackage[T1]{fontenc}
  \usepackage[utf8]{inputenc}
  \usepackage{textcomp} % provide euro and other symbols
\else % if luatex or xetex
  \usepackage{unicode-math} % this also loads fontspec
  \defaultfontfeatures{Scale=MatchLowercase}
  \defaultfontfeatures[\rmfamily]{Ligatures=TeX,Scale=1}
\fi
\usepackage{lmodern}
\ifPDFTeX\else
  % xetex/luatex font selection
\fi
% Use upquote if available, for straight quotes in verbatim environments
\IfFileExists{upquote.sty}{\usepackage{upquote}}{}
\IfFileExists{microtype.sty}{% use microtype if available
  \usepackage[]{microtype}
  \UseMicrotypeSet[protrusion]{basicmath} % disable protrusion for tt fonts
}{}
\makeatletter
\@ifundefined{KOMAClassName}{% if non-KOMA class
  \IfFileExists{parskip.sty}{%
    \usepackage{parskip}
  }{% else
    \setlength{\parindent}{0pt}
    \setlength{\parskip}{6pt plus 2pt minus 1pt}}
}{% if KOMA class
  \KOMAoptions{parskip=half}}
\makeatother
\ifLuaTeX
\usepackage[bidi=basic,shorthands=off]{babel}
\else
\usepackage[bidi=default,shorthands=off]{babel}
\fi
\ifLuaTeX
  \usepackage{selnolig} % disable illegal ligatures
\fi
\setlength{\emergencystretch}{3em} % prevent overfull lines
\providecommand{\tightlist}{%
  \setlength{\itemsep}{0pt}\setlength{\parskip}{0pt}}
\usepackage{bookmark}
\IfFileExists{xurl.sty}{\usepackage{xurl}}{} % add URL line breaks if available
\urlstyle{same}
\hypersetup{
  pdftitle={From bases to projections},
  pdflang={en},
  hidelinks,
  pdfcreator={LaTeX via pandoc}}

\title{From bases to projections}
\author{}
\date{}

\begin{document}
\maketitle

\section{From bases to projections}\label{from-bases-to-projections}

This foundation bridge supplies the finite-dimensional linear algebra
used to begin the averaging proof of complete reducibility. It belongs
to the programme's existing linear-algebra foundation; it is not a
replacement for that course.

The basis arguments below are adapted from Jim Hefferon's \emph{Linear
Algebra}, at source revision
\texttt{df2262e089a02651c127f1dd12649c4622ee1383}. The complete selected
source statements and proofs, their source locators, and the author's
grant are supplied with this lesson. This adaptation retains
\textbf{Creative Commons Attribution--ShareAlike 2.5} terms; it is not
CC0 and is not relicensed as GFDL. See the
\href{HEFFERON-LICENSE.txt}{author's grant},
\href{HEFFERON-ACKNOWLEDGEMENTS.txt}{acknowledgements},
\href{HEFFERON-PROOFS.tex}{exact selected source proofs}, and
\href{https://creativecommons.org/licenses/by-sa/2.5/}{reuse terms}.

\emph{Adaptation, scope comparison, examples and solutions by OpenAI
GPT-6 Astra in Codex, Ultra reasoning effort, October 2026. Self-checked
by the AI writing this bridge; no human or independent AI review is
claimed. The original source remains unchanged. This lesson corrects an
already-recorded coefficient sign, makes the field condition explicit,
and expands the finite termination and projection arguments. These are
standard results, not claimed discoveries.}

\subsection{What to know first}\label{what-to-know-first}

We use field and vector-space axioms, finite sums, finite induction and
elementary finite-set counting. A \textbf{field} has distinct elements
\(0\) and \(1\); every nonzero scalar has an inverse. A vector space
\(V\) over a field \(k\) is an abelian additive group with scalar
multiplication satisfying the vector-space axioms.

For \(S\subseteq V\), its \textbf{span} consists of finite linear
combinations of elements of \(S\); the empty combination is zero. The
set is \textbf{linearly independent} when every relation involving
finitely many distinct elements has all coefficients zero. A
\textbf{basis} is a linearly independent spanning set. Throughout, \(V\)
has a finite basis. No topology, inner product, algebraic closure or
characteristic-zero assumption is used.

Linear combinations of linear combinations remain linear combinations,
by distributivity. Thus a span is a subspace. A subset of an independent
set is independent, since a relation in the subset is also one in the
larger set. Finally, coordinates in a basis are unique: subtracting two
expressions gives a relation in the basis and therefore zero differences
between all coefficients. We will use these elementary consequences
explicitly.

\subsection{Adding one independent
vector}\label{adding-one-independent-vector}

\textbf{Lemma.} If \(S\) is independent and \(v\notin S\), then
\(S\cup\{v\}\) is independent exactly when
\(v\notin\operatorname{span}(S)\).

\textbf{Proof.} If \(v=\sum_{i=1}^r a_i s_i\), combine repeated terms so
the \(s_i\) are distinct. The equation

\[
\sum_{i=1}^r a_i s_i-v=0
\]

is a nontrivial relation on distinct elements of \(S\cup\{v\}\): the
coefficient of \(v\) is \(-1\ne0\). Conversely, any nontrivial relation
on that union must involve \(v\), because \(S\) is independent. Write it
as

\[
\sum_{i=1}^r a_i s_i+bv=0.
\]

If \(b=0\), independence of \(S\) forces all other coefficients to
vanish, a contradiction. Consequently \(b\ne0\), and

\[
v=-\sum_{i=1}^r b^{-1}a_i s_i\in\operatorname{span}(S).
\]

This proves both implications. It also covers \(S=\varnothing\): the
only dependent singleton is \(\{0\}\). \(\square\)

The minus sign in the last display repairs the existing source finding
\texttt{B40-VS2-SOURCE-003}. Omitting it did not change the source's
span-membership conclusion, but it made the displayed equality wrong
outside characteristic two.

\subsection{Exchanging a basis vector}\label{exchanging-a-basis-vector}

\textbf{Lemma.} Let \(B=(b_1,\ldots,b_n)\) be a basis, and write
\(v=\sum_{j=1}^n c_jb_j\). If \(c_i\ne0\), replacing \(b_i\) by \(v\)
gives another basis.

\textbf{Proof.} Suppose

\[
d_i v+\sum_{j\ne i}d_jb_j=0.
\]

Substitution expresses zero in the original basis. Its
\(b_i\)-coefficient is \(d_ic_i\), so \(d_i=0\). The other coefficients
are then \(d_j=0\). Thus the replacement family is independent.
Moreover,

\[
b_i=c_i^{-1}v-\sum_{j\ne i}c_i^{-1}c_jb_j.
\]

Every old basis vector therefore belongs to the span of the replacement
family. That family spans \(V\), and hence is a basis. \(\square\)

\subsection{Why finite dimension bounds
independence}\label{why-finite-dimension-bounds-independence}

\textbf{Theorem.} If \(V\) has a basis with \(n\) elements, every
independent subset of \(V\) has at most \(n\) elements. Consequently all
bases of \(V\) have the same size.

\textbf{Proof.} If \(n=0\), then \(V=\{0\}\), whose only independent
subset is empty. Suppose \(n>0\) and, towards a contradiction, that an
independent set contains distinct vectors \(s_1,\ldots,s_{n+1}\).

Start with the given basis \(B_0\). Inductively suppose that \(B_r\) is
a basis containing \(s_1,\ldots,s_r\) and \(n-r\) of the original basis
vectors, where \(0\le r<n\). Express \(s_{r+1}\) in \(B_r\). Some
coefficient of one of the remaining original vectors must be nonzero:
otherwise \(s_{r+1}\) would be in the span of \(s_1,\ldots,s_r\),
contradicting independence. The exchange lemma replaces such an original
vector by \(s_{r+1}\). This constructs \(B_{r+1}\).

After \(n\) exchanges, \((s_1,\ldots,s_n)\) is a basis. It spans
\(s_{n+1}\), again contradicting independence. Thus there cannot be
\(n+1\) distinct elements in an independent set. This proves the bound
without presupposing that the set spans or is already finite.

If another basis has \(m\) elements, the bound gives \(m\le n\).
Applying the same argument with that basis gives \(n\le m\). Hence
\(m=n\); their common value is \(\dim_k V\). \(\square\)

This is the source's exchange argument with its induction written
directly for an arbitrary independent set. It avoids using the existence
of a minimal-sized basis when proving the independent-set bound.

\subsection{Extending a basis and choosing a
complement}\label{extending-a-basis-and-choosing-a-complement}

\textbf{Theorem.} Every independent subset of \(V\) extends to a basis
of \(V\). Every subspace \(U\subseteq V\) has a finite basis, which
extends to one of \(V\).

\textbf{Proof.} An independent set \(S\) has at most \(n=\dim_k V\)
elements. If it does not span \(V\), choose a vector outside its span.
The first lemma says that adjoining it preserves independence. Repeat
whenever the span is still not \(V\). Each step increases the finite
cardinality by one. The preceding theorem makes more than \(n-|S|\)
additions impossible, so the process must reach a basis.

For a subspace \(U\), start with the empty set and choose each new
vector in \(U\) outside the current span. These sets are independent in
\(V\), so the same bound forces termination in a spanning independent
set of \(U\). Extend this basis to one of \(V\) by the first part. Only
finitely many choices are needed. If \(U=0\), its basis is empty; if
\(U=V\), there are no extra basis vectors. \(\square\)

\textbf{Corollary.} Every subspace of \(V\) is the image of a linear
projection. More precisely, there are linear maps \(i:U\to V\) and
\(p:V\to U\) with \(i\) the inclusion and \(pi=1_U\). The endomorphism
\(p_0=ip\) satisfies

\[
p_0^2=p_0,\qquad\operatorname{im}p_0=U,
\qquad V=U\oplus\ker p_0.
\]

\textbf{Proof.} Choose a basis \((u_1,\ldots,u_r)\) of \(U\) and extend
it to \((u_1,\ldots,u_r,w_1,\ldots,w_{n-r})\) of \(V\). Unique
coordinates define

\[
\begin{aligned}
v&=\sum_{i=1}^r a_i u_i+\sum_{j=1}^{n-r}b_jw_j,\\
p(v)&=\sum_{i=1}^r a_i u_i.
\end{aligned}
\]

Coordinates of a sum are the sums of the coordinates, and coordinates of
a scalar multiple are the scalar multiples of the coordinates, again by
uniqueness. Therefore \(p\) is linear. It fixes \(U\), which gives
\(pi=1_U\), \(p_0^2=p_0\), and \(\operatorname{im}p_0=U\). Its kernel is
the span of the \(w_j\). For every \(v\),

\[
v=p_0v+(v-p_0v),\qquad p_0(v-p_0v)=0.
\]

If \(u\in U\cap\ker p_0\), then \(u=p_0u=0\). This proves the direct
sum. \(\square\)

There is generally no distinguished choice of \(p_0\). In the displayed
adapted basis its matrix is \(\operatorname{diag}(I_r,0)\). A different
complement can give a different projection onto the same subspace.

\subsection{The exact step used in representation
theory}\label{the-exact-step-used-in-representation-theory}

In
\href{https://kokunoyumeto.github.io/open-mathematics-courses/courses/RT-FIN/representations-and-complete-reducibility.html}{Representations
and complete reducibility}, Theorem 2.3 begins by extending a basis of
an invariant subspace \(U\), then defining a projection \(p_0\). The
preceding corollary supplies exactly this construction, over every field
in that theorem, including finite fields. Exercise 5 applies the same
construction to its isotypic subspace \(S\). No inner product or
orthogonal projection is needed for either initial construction.

The representation-theoretic averaging has an additional condition:
\(|G|\) must be invertible in \(k\). If \(U\) is invariant, put

\[
P=\frac{1}{|G|}\sum_{g\in G}\rho(g)p_0\rho(g^{-1}).
\]

Here each summand maps into \(U\) and acts as the identity on \(U\).
Hence \(P(V)\subseteq U\) and \(P|_U=1_U\). Conjugation by \(\rho(h)\)
permutes the summands by \(g\mapsto hg\), so \(P\) commutes with the
group action. Also \(P^2=P\), because \(P\) fixes its image. The
decomposition \(v=Pv+(v-Pv)\) gives the invariant complement \(\ker P\).
The requirement that \(|G|\) be invertible belongs here, not to basis
extension.

This comparison applies to the exact lesson source SHA-256
\texttt{659942BDF47420E6665D59C4E80CF973867730D3B231D7E9C80AB8177BD83712},
at Theorem 2.3 and Exercise 5. It checks only the
finite-basis/projection prerequisite and its use. It does not certify
the rest of the lesson, its other prerequisites, the isotypic
identification, or the complete course. The original consumer lesson and
its review state are unchanged.

\subsection{Worked example and
exercises}\label{worked-example-and-exercises}

\textbf{Example.} In \(k^2\), take \(U=k(1,0)\). For every \(a\in k\),
the vector \((a,1)\) is a complement basis vector, and the resulting
projection is

\[
p_a(x,y)=(x-ay,0),\qquad
[p_a]=\begin{pmatrix}1&-a\\0&0\end{pmatrix}.
\]

Its kernel is \(k(a,1)\). Direct substitution gives \(p_a^2=p_a\). These
distinct projections show why choosing a subspace does not canonically
choose a complement.

\textbf{Exercise 1.} Over a field, extend \(u=(1,1,0)\), \(v=(0,1,1)\)
to a basis of \(k^3\), and write a projection onto their span.

\textbf{Solution.} A relation \(au+bv=0\) gives \(a=0\) in the first
coordinate and \(b=0\) in the third. The vector \(e_2=(0,1,0)\) is
outside their span: the same two coordinates would force both
coefficients to vanish. Thus \((u,v,e_2)\) is a basis. In fact

\[
\begin{aligned}
(x,y,z)&=xu+zv+(y-x-z)e_2,\\
p_0(x,y,z)&=(x,x+z,z).
\end{aligned}
\]

This formula is valid in every characteristic. It fixes \(u,v\), kills
\(e_2\), and satisfies \(p_0^2=p_0\).

\textbf{Exercise 2.} Why does basis extension not by itself prove that
every invariant subspace has an invariant complement in characteristic
\(p\)?

\textbf{Solution.} Over a field of characteristic \(p>0\), let a
generator of \(C_p\) act on \(k^2\) by

\[
A=\begin{pmatrix}1&1\\0&1\end{pmatrix}=I+N.
\]

Since \(N^2=0\), induction gives \((I+N)^m=I+mN\), and thus \(A^p=I\).
The line \(U=ke_1\) is invariant. Any complementary line has a generator
\(w=(a,1)\). If it were invariant, \(Aw=(a+1,1)\) would equal \(cw\).
The second coordinate forces \(c=1\), while the first then forces
\(1=0\), impossible. Ordinary linear complements exist; invariant ones
do not in this example. Averaging cannot repair them because the scalar
\(|C_p|=p\) is zero in \(k\).

\textbf{Exercise 3.} Which part of the proof prevents it from asserting
basis extension for arbitrary infinite-dimensional spaces?

\textbf{Solution.} The finite bound on the size of an independent set
makes the process of adding vectors terminate. With no finite spanning
basis, this argument supplies no such bound and therefore no termination
claim. An infinite-dimensional theorem needs a different argument and
any additional set-theoretic assumptions must be stated. Nothing in this
bridge claims that theorem.

\subsection{Sources and changes}\label{sources-and-changes}

Hefferon, \emph{Linear Algebra}, vector-spaces chapter: the add-vector
lemma in \texttt{src/vs/vs2.tex}, lines 492--525; exchange, equal basis
size, independent-set bound and basis extension in
\texttt{src/vs/vs3.tex}, lines 1657--1779, 1818--1831 and 1857--1874.
The
\href{https://kokunoyumeto.github.io/program-matematika-indonesia/en/readers/hefferon-foundations/}{programme's
existing foundation reader} preserves the source edition and its
editorial notes. The \href{https://hefferon.net/linearalgebra/}{author's
site} supplies the broader course.

The adaptation above keeps the same algebraic argument, spells out the
zero-dimensional case and finite termination, proves the subspace and
projection consequences, and gives examples with solutions. The
requirement \(0\ne1\) also respects the already-recorded source-context
finding \texttt{B40-FIELDS-SOURCE-001}. It is not a newly discovered
defect. The field-general formulation is justified by the displayed
algebraic proof, not merely by changing the word ``real'' in a source
statement.

All adapted teaching and added exposition on this page use
\href{https://creativecommons.org/licenses/by-sa/2.5/}{CC BY-SA 2.5}.
The separately linked representation-theory lesson retains its own
source, credit and licence. No change to it or to the canonical Stacks
text is made here.

\end{document}
